\documentclass[tikz,border=10pt]{standalone}
\usepackage{ctex} % 支持中文显示
\usepackage{xcolor}

% fig12 专属配色：四列设备分色
\definecolor{colSend}{RGB}{41,98,255}    % 发送端 蓝
\definecolor{colSendBg}{RGB}{230,238,255}
\definecolor{colToR}{RGB}{240,140,0}      % ToR (HoCM) 橙 高亮
\definecolor{colToRBg}{RGB}{255,243,224}
\definecolor{colCore}{RGB}{120,120,120}   % 核心交换机 灰
\definecolor{colCoreBg}{RGB}{242,242,242}
\definecolor{colRecv}{RGB}{46,160,87}     % 接收端 绿
\definecolor{colRecvBg}{RGB}{228,245,233}

\begin{document}
\begin{tikzpicture}[>=stealth, thick]

% 样式定义
\tikzstyle{innerbox} = [draw, solid, thick, rectangle, minimum width=2.2cm, minimum height=1.2cm, align=center, font=\songti, fill=white]
\tikzstyle{title} = [align=center, font=\songti\bfseries]

% 定义各列的中心 X 坐标
\def\xS{0}     % 发送端主机
\def\xT{4.5}   % ToR交换机
\def\xC{9}     % 核心交换机
\def\xR{13.5}  % 接收端主机

% 绘制四个外部大方框 (实线) —— 各列分色 + 浅色底
\def\yTop{2.5}
\def\yBot{-7.5}
\def\hw{1.6} % 外部大框的宽度一半

\filldraw[draw=colSend, fill=colSendBg, thick] (\xS-\hw, \yTop) rectangle (\xS+\hw, \yBot);
\filldraw[draw=colToR,  fill=colToRBg,  very thick] (\xT-\hw, \yTop) rectangle (\xT+\hw, \yBot);
\filldraw[draw=colCore, fill=colCoreBg, thick] (\xC-\hw, \yTop) rectangle (\xC+\hw, \yBot);
\filldraw[draw=colRecv, fill=colRecvBg, thick] (\xR-\hw, \yTop) rectangle (\xR+\hw, \yBot);

% 顶部标题
\node[title, text=colSend!80!black] at (\xS, 1.6) {发送端主机};
\node[title, text=colToR!70!black] at (\xT, 1.6) {ToR交换机 \\ (HoCM启用)};
\node[title, text=colCore!60!black] at (\xC, 1.6) {核心交换机 \\ (标准转发)};
\node[title, text=colRecv!70!black] at (\xR, 1.6) {接收端主机};

% 发送端主机内部节点
\node[innerbox, draw=colSend] (S1) at (\xS, 0) {DMA队列 \\ 采样};
\node[innerbox, draw=colSend] (S2) at (\xS, -2) {CL/BT \\ 嵌入};

% ToR交换机内部节点
\node[innerbox, draw=colToR] (T1) at (\xT, 0) {Parser \\ 提取CL/BT};
\node[innerbox, draw=colToR] (T2) at (\xT, -2) {FSM更新};
\node[innerbox, draw=colToR] (T3) at (\xT, -4) {调度决策};
\node[innerbox, draw=colToR] (T4) at (\xT, -6) {队列调度};

% 核心交换机内部节点
\node[innerbox, draw=colCore] (C1) at (\xC, 0) {Parser \\ 标准转发};

% 接收端主机内部节点
\node[innerbox, draw=colRecv] (R1) at (\xR, 0) {应用层 \\ 处理};

% 横向数据流箭头连接
\draw[->, colSend] (S1.east) -- (T1.west);
\draw[->, colToR] (T1.east) -- (C1.west);
\draw[->, colCore] (C1.east) -- (R1.west);

\end{tikzpicture}
\end{document}
